\section{SF-E844. A non-power complete-intersection family satisfying the socle-pair criterion}\label{SF-E844:sf-e844.-a-non-power-complete-intersection-family-satisfying-the-socle-pair-criterion}

Author: /root/sixfolds, GPT-6 Astra Pro; inherited effort not exposed.
Date: 2026-11-22.

\textbf{Status:} the complete author, root and separate independent peer whole-proof readings have passed. This is an explicit derived application of the already accepted SF-E842 criterion, not a separate reduction of the residual gate defined by that criterion.
The theorem SF-E842 is an already whole-reviewed input; this note does not
modify that frozen theorem or claim its socle criterion for every complete
intersection. The new point is an explicit family outside the parameter-power
scope of E843 and all the previously listed scalar/commuting-triple/quadratic
tests. No semiregular seed is constructed.

\subsection{1. The actual complete intersections and their finite algebra}\label{SF-E844:the-actual-complete-intersections-and-their-finite-algebra}

Use the same constant Poisson coordinates and integral mixing matrix as E843:
\(X=(x_1,x_2,y_1,y_2)\), with
\(\pi_+=\partial_{x_1}\wedge\partial_{x_2}\) and
\(\pi_-=\partial_{y_1}\wedge\partial_{y_2}\), and put \(u=JX\), where

\[
 J=\begin{pmatrix}
 1&0&1&0\\
 0&1&0&1\\
 -3&5&-9&10\\
 -8&5&-16&15
 \end{pmatrix},\qquad \det J=-20.
 \tag{E844.1}
\]

For any nonzero complex number \(\lambda\), define

\[
 f_1=u_1^3+\lambda u_2^2u_3^2,\quad
 f_2=u_2^3,\quad f_3=u_3^3,\quad f_4=u_4^3,\qquad
 I_\lambda=(f_1,f_2,f_3,f_4),\quad
 A_\lambda=\mathbb C[[u_1,u_2,u_3,u_4]]/I_\lambda.
 \tag{E844.2}
\]

The algebra is finite and is a complete intersection. To verify finiteness
and its exact basis without a flatness assumption, set
\(B=\mathbb C[u_2,u_3,u_4]/(u_2^3,u_3^3,u_4^3)\).
The polynomial quotient
\(B[u_1]/(u_1^3+\lambda u_2^2u_3^2)\) is free over \(B\), with basis
\(1,u_1,u_1^2\), because the polynomial in \(u_1\) is monic.
Its generators are all nilpotent:
\(u_1^6=\lambda^2u_2^4u_3^4=0\) and \(u_j^3=0\) for \(j=2,3,4\).
It therefore equals its formal completion and is \(A_\lambda\).
Consequently

\[
 \{u_1^{r_1}u_2^{r_2}u_3^{r_3}u_4^{r_4}:0\le r_i\le2\}
 \text{ is a basis},\qquad \dim_{\mathbb C}A_\lambda=81.
 \tag{E844.3}
\]

The ideal is generated by four elements and has height four in the regular
local ring of dimension four; hence they form a regular sequence.
Alternatively, the first three coordinate powers \(u_2^3,u_3^3,u_4^3\)
form a regular sequence and the remaining monic relation is a nonzerodivisor,
which can be checked by successive powers of the nilpotent ideal
\((u_2,u_3,u_4)\).

The only nonzero reduction of a forbidden exponent is

\[
                   u_1^3=-\lambda u_2^2u_3^2.
 \tag{E844.4}
\]

This reduction increases total degree by one. A monomial of total degree
at least nine either vanishes during reduction or would have to reduce
to a basis monomial of degree at least nine, whereas all basis monomials
have degree at most eight. Thus the ninth maximal-ideal power is zero.
The degree filtration of the displayed basis has associated graded algebra
\(\mathbb C[u_1,u_2,u_3,u_4]/(u_1^3,u_2^3,u_3^3,u_4^3)\).
In particular, the lowest nonzero ideal order is three and the embedding
dimension is four.

The element
\(\tau=u_1^2u_2^2u_3^2u_4^2\) is nonzero by the basis and is killed by
every \(u_i\), including \(u_1\) by (E844.4).
The complete-intersection Jacobian-socle theorem, in its precise
characteristic-zero form already used in E842, gives
\(\operatorname{Soc}(A_\lambda)=\mathbb C\tau\). The exact published locator is Eisenbud--Ulrich, Theorem 7.1, p.~210, with Appendix A.1/A.3/A.5, pp.~222--225.\cite{SFEisenbudUlrich2019}

\subsection{2. The two actual bracket entries form a socle pair}\label{SF-E844:the-two-actual-bracket-entries-form-a-socle-pair}

The constant coefficient tuples of \(\pi_+,\pi_-\) in the \(u\)-coordinates,
in the order \(12,13,14,23,24,34\), are

\[
 p=(1,5,5,3,8,25),\qquad q=(1,10,15,9,16,25).
 \tag{E844.5}
\]

Put \(P_\pm=(\pi_\pm(df_i,df_j))_{i<j}\) in \(A_\lambda^6\), and
\(P=P_++P_-\). The derivative of the additional term in \(f_1\) is

\[
 d(\lambda u_2^2u_3^2)
       =2\lambda u_2u_3^2\,du_2+2\lambda u_2^2u_3\,du_3.
 \tag{E844.6}
\]

In the bracket with \(f_2=u_2^3\), the first term pairs trivially with
\(du_2\), while the second contains \(u_2^4\) and vanishes.
The unmodified pair \(f_3,f_4\) is immediate. Therefore the two exact
pencil entries, including all higher terms, are

\[
             P_{12}=18u_1^2u_2^2,\qquad
             P_{34}=450u_3^2u_4^2.
 \tag{E844.7}
\]

Let \(E_{12}=u_1^2u_2^2\), \(E_{34}=u_3^2u_4^2\).
The ideal \(E_{12}A_\lambda\) has the nine basis monomials
\(E_{12}u_3^au_4^b\), \(0\le a,b\le2\).
Indeed multiplication by \(u_2\) kills it, and multiplication by \(u_1\)
also kills it because \(u_1E_{12}=-\lambda u_2^4u_3^2=0\).
Likewise, \(E_{34}A_\lambda\) has basis
\(E_{34}u_1^au_2^b\), \(0\le a,b\le2\).
The relation \(u_1^3=-\lambda u_2^2u_3^2\) contributes nothing after
multiplication by \(E_{34}\), because it then contains \(u_3^4\).
The two indicated basis subsets meet exactly in \(\tau\). Hence

\[
 P_{12}A_\lambda\cap P_{34}A_\lambda=\mathbb C\tau,\qquad
                        P_{12}P_{34}=8100\tau\ne0.
 \tag{E844.8}
\]

This is precisely the sufficient condition in SF-E842. No comparison with a
different quotient, and no base change of a deformation, has been used.

The Jacobian provides a further exact check. In the \(u\)-coordinates the
derivative matrix is triangular apart from the two entries in the first row,
so its determinant is \(81\tau\). Returning to \(X\) gives

\[
 \det\bigl(\partial f_i/\partial X_j\bigr)
       =81\det(J)\tau=-1620\tau,\qquad
                    P_+\wedge P_-=-1620\tau.
 \tag{E844.9}
\]

The wedge convention is the one fixed in E843.4.

\subsection{3. These are not parameter-power ideals in any formal coordinates}\label{SF-E844:these-are-not-parameter-power-ideals-in-any-formal-coordinates}

Suppose, for contradiction, that there are formal parameters \(v_1,\ldots,v_4\)
and positive integers \(a_i\) such that
\(I_\lambda=(v_1^{a_1},\ldots,v_4^{a_4})\).
There is no element of ideal order less than three, so every \(a_i\ge3\).
The degree-three initial relation space of \(I_\lambda\) has dimension four.
Only generators with \(a_i=3\) can contribute to this space, so all four
exponents must equal three.

Write \(\ell_i\) for the linear part of \(v_i\).
Since \(v_i^3\in I_\lambda\), its initial cubic belongs to

\[
                 \operatorname{span}_{\mathbb C}
                           (u_1^3,u_2^3,u_3^3,u_4^3).
 \tag{E844.10}
\]

A nonzero linear form whose cube belongs to this span has only one nonzero
coordinate: otherwise its cube has a nonzero coefficient \(3a^2b\)
on a mixed monomial \(u_j^2u_k\).
The four \(\ell_i\) are independent, so after reordering and rescaling the
\(v_i\) we may assume \(\ell_i=u_i\).

Write \(v_1=u_1+Q_1+O(\mathfrak m^3)\), with \(Q_1\) homogeneous quadratic.
Because \(v_1^3\) belongs to the ideal generated by the four \(f_j\), write
\(v_1^3=\sum_j a_{1j}(u)f_j\).
Comparison of degree-three terms gives \(a_{11}(0)=1\) and
\(a_{1j}(0)=0\) for \(j\ne1\).
Comparison of degree-four terms therefore gives, for certain linear forms
\(L_j\),

\[
              3u_1^2Q_1
                 =\lambda u_2^2u_3^2+\sum_{j=1}^4L_j u_j^3.
 \tag{E844.11}
\]

This is impossible. The monomial \(u_2^2u_3^2\) is divisible neither by
\(u_1^2\) nor by any \(u_j^3\), so its coefficient is zero on the left
and \(\lambda\ne0\) on the right. This excludes every formal parameter
change, without assuming that it preserves the Poisson planes.

\subsection{4. The old local tests still do not exclude this family}\label{SF-E844:the-old-local-tests-still-do-not-exclude-this-family}

Every bracket has order at least four. The additional terms from (E844.6)
have order at least five. The product of a higher bracket term with any
bracket term has order at least nine and is therefore zero in \(A_\lambda\).
Thus every scalar cross-minor is determined by the original quartic terms:

\[
                P_{+,e}P_{-,f}-P_{-,e}P_{+,f}=0
                       \quad\text{for every ordered }e,f.
 \tag{E844.12}
\]

For a direct verification, if the two edge indices overlap outside
coordinate \(1\), the product contains \(u_j^4\) with \(j\ge2\) and vanishes.
If the shared index is \(1\), replacing \(u_1^3\) by (E844.4) raises the
degree of the degree-eight product to at least nine, so it vanishes too.
For complementary edges the product is \(\tau\); the coefficient vanishes
because the ratios \(q_e/p_e\) agree on each complementary pair, as in E843.
This proves all scalar cross-minors, rather than only selected minors.

The initial pure bracket maps have rank six, since all coefficients in
(E844.5) are nonzero and the six monomials \(u_i^2u_j^2\) are independent
in the degree-four associated graded quotient.
By the exact covariance under arbitrary minimal-generator changes,
no independent generator triple can commute for either pure bivector.
Thus E836's commuting-triple criterion does not apply in any presentation.

The ideal has lowest order three and embedding dimension four, so it is
outside the square and regular quadratic-jet families E837--840, as well as
the lower-embedding-dimension part of E836. Section 3 excludes E843's full
parameter-power family, even allowing nonlinear mixed coordinates.
The assertion here is precisely an explicit additional application of E842;
it is not a claim that any unclassified complete intersection satisfies E842.

\subsection{5. Actual coherent objects excluded by the criterion}\label{SF-E844:actual-coherent-objects-excluded-by-the-criterion}

On the genuine RM abelian fourfold with the reviewed balanced locally free
hull \(H\), realize this quotient at a point as a finite subscheme \(Z_\lambda\).
This is algebraic because \(I_\lambda\) contains the ninth maximal-ideal
power. Take sufficiently many copies of \(H\), choose a surjection
\(H\to\mathcal O_{Z_\lambda}^m\), and let \(K\) be its kernel.
Every extension

\[
 0\longrightarrow\mathcal O_{Z_\lambda}^m
   \longrightarrow E\longrightarrow K\longrightarrow0
 \tag{E844.13}
\]

is an actual coherent, hence perfect and ordinarily gluable, sheaf.
Its torsion is exactly \(\mathcal O_{Z_\lambda}^m\), its double dual is \(H\),
its hull quotient is \(\mathcal O_{Z_\lambda}^m\), and
\(\operatorname{ch}(E)=\operatorname{ch}(H)\).
The same exact sequence and complete-intersection Koszul calculation as in
E843 give the actual extension-space dimension \(486m^2\).

For any specified allowed coherent effective finite translation/tensor action,
weak balanced middle semiregularity gives the two local ordinary bivector
deformations by the reviewed E819 interface, also with a rational algebraic
\(B\)-field. The socle pair (E844.8) and E842 exclude those deformations.
Thus every extension (E844.13) fails the required semiregularity.
No equivariant splitting or simplicity is assumed.

\subsection{6. Exact certificate and review boundary}\label{SF-E844:exact-certificate-and-review-boundary}

The companion criterion and universal all-multiplicity obstruction are the
written E842 proof. A separate exact quotient-ring program reconstructs this
family, its multiplication and actual gradients; its numerical scope is
recorded in the adjacent JSON. The non-power assertion is the formal
degree-three/four argument in Section 3, not an inference from finitely many
coordinate changes. The separate independent peer has read the entire draft and canonical proof, the whole author checker and its complete successful output, and has executed a distinct exact symbolic reconstruction. Root read the entire draft, canonical proof and peer report and compared all thirteen displays. No fresh root code execution or whole-code reading is claimed for this example.

\paragraph{Certificate provenance.}
The exact BigInt checker passed at \texttt{2026-11-22T01:17:56.329Z} in \(0.060296346\) seconds. It checks the \(\lambda=1\) witness, including all \(6561\) basis products, \(531441\) associativity triples, \(220\) degree-nine monomials, \(36\) scalar cross-minors, the actual Jacobian, and both principal-ideal dimensions. The first attempt stopped with a Boolean--BigInt precedence error; the original source, empty standard output, error and one-parenthesis repair are preserved. No mathematical revision was required. The universal \(\lambda\ne0\), formal-coordinate and all-multiplicity assertions rest on the written proof.
